\section{Related Works}
\paragraph{Large Language Model Compression Techniques.}
The rapid growth of LLMs has spurred extensive research on compression methods that reduce model size and inference cost while preserving performance. Existing approaches can be broadly categorized into four classes: pruning, quantization~\citep{Nagel2021WhitePaperQuantization,frantar2022gptq,hu2025ostquant,lin2023awq,Hu2024I-LLM}, knowledge distillation~\citep{hinton2015distilling}, and low-rank approximation~\citep{yuan2023asvd,Wang2024SVD-LLM}. Unstructured pruning~\citep{ashkboos2024slicegpt,zhang2023pruning} removes individual weights, but its irregular sparsity patterns often yield limited hardware benefits. Structured pruning~\citep{ma2023llmpruner,dettmers2023spqr} eliminates entire channels or blocks, improving efficiency but risking significant accuracy loss. Quantization~\citep{frantar2022gptq,lin2023awq,hu2025ostquant} reduces numerical precision, offering strong memory savings but typically restricted to 2–8 bits, which limits flexibility and may degrade performance under aggressive settings. Knowledge distillation transfers knowledge from a large teacher model to a smaller student~\citep{liu-etal-2024-llm}, but requires costly retraining. In contrast, low-rank approximation, especially singular value decomposition (SVD), provides a hardware-agnostic, post-training alternative that can be combined with other methods, making it a particularly attractive solution for LLM compression.

\paragraph{SVD-based Compression for LLMs.}
SVD~\citep{Golub1987Generalization} reduces matrix sizes by truncating the smallest singular values. 
Low-rank approximation using SVD has been widely studied as a hardware-agnostic and efficient approach for compressing LLMs.
%
Early works mainly focused on minimizing the per-layer truncation loss.
ASVD~\citep{yuan2023asvd} introduced a learnable scaling matrix to mitigate the loss of discarded singular values, but the improvement remained local and could not guarantee the theoretical minimum truncation loss across all layers. SVD-LLM~\citep{Wang2024SVD-LLM,Wang2025SVD-LLM-V2} addressed this limitation by incorporating a whitening matrix that normalizes input activations, thereby achieving theoretically minimal truncation error. AdaSVD~\citep{li2025adasvd} introduces an alternating optimization scheme that iteratively updates low-rank matrices and compensation terms, thereby progressively reducing compression error. Several subsequent methods have incorporated importance weighting. For example, FW-SVD~\citep{Hsu2022FWSVD} and GFW-SVD~\citep{Chekalina2025GFWSVD} leverage global importance indicators to guide single-layer compression, whereas DipSVD~\citep{Ding2025DipSVD} jointly exploits both local and global importance to refine the weighting process.
%
Although these techniques improve layer-level stability, they still rely on the assumption that layer errors are independent, leaving the problem of cumulative error unresolved.
